\section{Worst-Case Bias Comparison}\label{sec:bias}

This section compares the asymptotic worst-case bias of the pooled and segmented estimators.
The main result shows that segmentation reduces worst-case bias whenever the conditional density $f_{Z|X}(z|x)$ varies with $x$.
Unlike the discrete-covariate case, where stratification fully eliminates pooling bias, segmentation with segments of positive width $\delta > 0$ only \emph{reduces} the compositional term --- within each segment, the conditional density can still vary with $x$, generating residual pooling bias.

Throughout, I consider the local constant estimator ($q = 0$) with a uniform kernel and equal segment weights $w_j = 1/S$.


\subsection{Worst-Case Bias of the Pooled Estimator}\label{sec:bias:pooled}

The base RD estimator \eqref{eq:base_rd} targets the pooled regression function $\bar{m}(x) = \int m(x,z)\, f_{Z|X}(z|x)\, dz$.
By \textcite{armstrong_simple_2020}, the worst-case bias of the local constant estimator with bandwidth $h$ and uniform kernel over a Lipschitz class with constant $M_P$ is
\begin{equation}\label{eq:bias_pooled}
	\wcb(\hat{\tau}(h)) = h\, M_P\, B_{1,0}
\end{equation}
where $B_{1,0} = \frac{1}{2}$ is the kernel constant for the local constant estimator with uniform kernel, and $M_P$ bounds $\abs{\bar{m}'(x)}$ and $\abs{\bar{m}'(x) + \bar{\tau}'(x)}$.

The bound $M_P$ must be chosen to be valid for the pooled regression function.
From the derivative decomposition \eqref{eq:mbar_deriv},
\begin{equation}\label{eq:MP_decomp}
	\bar{m}'(x) = \int \frac{\partial m}{\partial x}(x,z)\, f_{Z|X}(z|x)\, dz + \int m(x,z)\, \frac{\partial f_{Z|X}}{\partial x}(z|x)\, dz.
\end{equation}
The first term is bounded by $M_x$ (the Lipschitz constant of $m$ in $x$).
The second term reflects compositional changes and depends on the heterogeneity of $m(x,z)$ across $z$ and the rate of change of $f_{Z|X}$ in $x$.

Applying the triangle inequality:
\begin{equation}\label{eq:MP_bound}
	M_P \leq M_x + \sup_x \abs*{\int m(x,z)\, \frac{\partial f_{Z|X}}{\partial x}(z|x)\, dz} \eqdef M_x + \Delta_{\mathrm{comp}}.
\end{equation}
The term $\Delta_{\mathrm{comp}}$ is the full compositional bias contribution.
It vanishes if and only if either the regression function is constant in $z$ (no heterogeneity) or $f_{Z|X}(z|x)$ does not depend on $x$ (no compositional changes).


\subsection{Worst-Case Bias of the Segmented Estimator}\label{sec:bias:segmented}

The segmented estimator with equal weights is $\hat{\tau}_S(h) = \frac{1}{S} \sum_{j=1}^{S} \hat{\tau}_j(h)$.
Within segment $S_j$, the estimator $\hat{\tau}_j(h)$ is a local constant RD estimator applied to the subsample $\{i : Z_i \in S_j\}$.
This estimator implicitly targets the within-segment pooled regression function
\begin{equation}\label{eq:m_j}
	\bar{m}_j(x) \defeq \frac{\int_{S_j} m(x,z)\, f_{Z|X}(z|x)\, dz}{\int_{S_j} f_{Z|X}(z|x)\, dz}.
\end{equation}
Unlike the discrete case, where stratification fully conditions on the covariate value, the within-segment function $\bar{m}_j(x)$ still averages over a range of $z$-values of width $\delta$.
Its derivative inherits a within-segment compositional term:
\begin{equation}\label{eq:m_j_deriv}
	\bar{m}_j'(x) = \underbrace{\frac{\int_{S_j} \frac{\partial m}{\partial x}(x,z)\, f_{Z|X}(z|x)\, dz}{\int_{S_j} f_{Z|X}(z|x)\, dz}}_{\text{within-segment smoothing}} + \underbrace{\text{within-segment compositional term}}_{\text{reduced but not eliminated}}.
\end{equation}
The first term is bounded by $M_x$.
The within-segment compositional term depends on how much $m(x,z)$ varies across $S_j$ in the $z$-direction, which is at most $M_z \cdot \delta$ under the anisotropic class $\mathcal{F}_1^{\mathrm{aniso}}(M_x, M_z)$.
Consequently, the within-segment compositional contribution scales with $\delta$:
\begin{equation}\label{eq:delta_comp_j}
	\abs{\bar{m}_j'(x)} \leq M_x + \Delta_{\mathrm{comp},j}(\delta),
\end{equation}
where $\Delta_{\mathrm{comp},j}(\delta) \to 0$ as $\delta \to 0$, with the rate depending on $M_z$ and the density gradient.


\subsection{Bias Comparison}\label{sec:bias:comparison}

\begin{proposition}[Worst-Case Bias Reduction from Segmentation]\label{prop:bias_comparison}
	Assume $m(x,z) \in \mathcal{F}_1^{\mathrm{aniso}}(M_x, M_z)$.
	Then for any bandwidth $h > 0$ and segment width $\delta > 0$,
	\begin{equation}\label{eq:bias_comparison}
		\wcb(\hat{\tau}_S(h)) \leq \wcb(\hat{\tau}(h)),
	\end{equation}
	with strict inequality whenever $\Delta_{\mathrm{comp}} > \Delta_{\mathrm{comp},j}(\delta)$ for some $j$, i.e., whenever segmentation reduces the compositional term relative to full pooling.
\end{proposition}

The magnitude of the bias reduction depends on the segment width $\delta$.
With $\delta = h$, the within-segment compositional term is of order $M_z \cdot h$, making the residual pooling bias of order $h^2 \cdot M_z$ --- which is higher order relative to the leading smoothing bias term $h \cdot M_x$.
Full pooling, by contrast, has a compositional term of order $h \cdot \Delta_{\mathrm{comp}}$ that enters the leading bias.


\subsection{Worst-Case Function}\label{sec:bias:worst_case}

To build intuition for the source and magnitude of pooling bias, consider the worst-case function that maximizes the compositional term.

\paragraph{Isolating pooling bias.}
Setting $\frac{\partial m}{\partial x}(x,z) = 0$ eliminates the standard smoothing bias and isolates the compositional component.
Under this restriction, $m(x,z) = m(z)$ depends only on the covariate, and the pooled regression function becomes
\begin{equation}\label{eq:mbar_pure_comp}
	\bar{m}(x) = \int m(z)\, f_{Z|X}(z|x)\, dz.
\end{equation}
Despite $m$ being constant in $x$, the pooled function $\bar{m}(x)$ varies with $x$ through the changing density $f_{Z|X}(z|x)$.
Its derivative is
\begin{equation}\label{eq:mbar_prime_pure}
	\bar{m}'(x) = \int m(z)\, \frac{\partial f_{Z|X}}{\partial x}(z|x)\, dz = \Cov_{Z|X}\!\left(m(Z),\; \frac{\partial \log f_{Z|X}(Z|x)}{\partial x} \;\bigg|\; X = x\right),
\end{equation}
which is the covariance between the regression function and the score of the conditional density.
This quantity is large when $m(z)$ and $\frac{\partial f_{Z|X}}{\partial x}$ are strongly associated.

\paragraph{Maximizing heterogeneity.}
Under the anisotropic smoothness class, $m(z)$ is Lipschitz in $z$ with constant $M_z$.
The worst case sets the $z$-derivative at its bound:
\begin{equation}\label{eq:worst_case_m}
	m^*(z) = M_z \cdot z,
\end{equation}
so that $m$ increases linearly and maximally in $z$.

\paragraph{Worst-case density.}
Given $m^*(z) = M_z \cdot z$, the compositional term $\bar{m}'(x) = M_z \int z\, \frac{\partial f_{Z|X}}{\partial x}(z|x)\, dz$ is maximized when $\frac{\partial f_{Z|X}}{\partial x}(z|x)$ is positively associated with $z$: the density should shift mass toward high $z$ as $x$ increases (or equivalently, shift mass away from high $z$ as $x$ decreases).

Concretely, the worst-case density concentrates mass at high values of $z$ on one side of the cutoff and shifts this mass toward low $z$ on the other side.
In the $(x,z)$-plane, this corresponds to a density that is highest in the upper-left region (high $z$, low $x$) and declines toward the lower-right region (low $z$, high $x$), with the steepest possible gradient.

\paragraph{Geographic interpretation.}
This worst case has a natural geographic analog.
Suppose a city is located along the treatment boundary at a specific position (high $z$).
Close to the boundary (small $|x|$), most observations come from the city, so $f_{Z|X}(z|x)$ is concentrated near the city's location.
Further from the boundary (larger $|x|$), the city's influence wanes and the population is more dispersed or concentrated elsewhere.
If outcomes $m(z)$ differ between urban and rural locations, pooling across the boundary mixes these compositional shifts into the estimated treatment effect, generating pooling bias.
Segmentation mitigates this by restricting comparisons to units at similar positions along the boundary.

\paragraph{Effect of segmentation on the worst case.}
Under segmentation with segment width $\delta$, the worst-case function $m^*(z) = M_z \cdot z$ can vary by at most $M_z \cdot \delta$ within each segment.
The within-segment compositional term is therefore bounded by a quantity proportional to $M_z \cdot \delta$ times the within-segment density gradient.
As $\delta \to 0$, this term vanishes --- but for any finite $\delta > 0$, residual pooling bias remains.
The practical question is whether $\delta = h$ makes the residual term small enough to be higher-order relative to the leading smoothing bias.
